\documentclass{article}

% NeurIPS 2026 main-track camera-ready version. The style file defines the
% camera-ready footer only when the track option (main) is passed with final.
\PassOptionsToPackage{numbers,sort&compress}{natbib}
\usepackage[main,final]{neurips_2026}

\usepackage[utf8]{inputenc}
\usepackage[T1]{fontenc}
\usepackage{hyperref}
\usepackage{url}
\usepackage{booktabs}
\usepackage{amsmath}
\usepackage{amsfonts}
\usepackage{nicefrac}
\usepackage{microtype}
\usepackage{enumitem}
\usepackage{xcolor}
\usepackage{array}
\usepackage{graphicx}
\usepackage{tikz}
\usetikzlibrary{arrows.meta,calc,positioning,shadows}

\hypersetup{
  pdftitle={Measuring Safety Alignment Effects in Autonomous Security Agents},
  pdfauthor={Isaac David, Arthur Gervais}
}
\newcommand{\codeurl}{\url{https://github.com/zaddy6/measuring-safety-alignment-effects}}

\title{Measuring Safety Alignment Effects in Autonomous Security Agents}
\author{%
  Isaac David \\
  University College London \\
  \texttt{isaac.david.22@ucl.ac.uk} \\
  \And
  Arthur Gervais \\
  University College London \\
  \texttt{a.gervais@ucl.ac.uk} \\
}
\date{}

\begin{document}
\maketitle

\begin{abstract}
Autonomous security agents use language models to inspect code, call tools, and test vulnerabilities inside authorized environments. Most safety evaluations ask whether a model refuses harmful single-turn requests. We study a different failure mode: whether safety alignment changes the evidence-grounded behavior of an agent that is already operating in a local sandbox with fixed tools and success checks. We evaluate four pairs of safety-aligned models and less-restricted (uncensored or abliterated) derivatives in the same security-agent harness: Gemma 4 31B, Gemma 4 26B A4B, Qwen2.5-Coder 7B, and Llama 3.1 8B. Each model receives authorized vulnerability-analysis tasks, and each run is scored from saved traces rather than self-reported success. We measure task completion, refusals, unsafe actions, and whether the final artifact is grounded in the relevant files, symbols, and vulnerability evidence. The key finding is that the gap is not mainly visible refusal. The aligned Gemma 4 models usually continue working, but their artifacts are less likely to satisfy security-specific evidence checks. The uncensored Gemma 4 derivatives more often find the relevant code, identify reachability, and write usable security reports. Clear authorization in the prompt does not recover this behavior for Gemma 4 31B. The pattern is not general: the Gemma derivatives also do better on non-security coding controls, and the abliterated Qwen2.5-Coder and Llama 3.1 derivatives do not outperform their aligned counterparts. At the same time, all conditions fail the hardest proof-of-trigger and patch-verification tasks. These results suggest that safety alignment effects in autonomous security agents cannot be understood only as refusals: they also appear as weaker grounding and less specific defensive evidence, and they depend on model family, derivative, and tool interface. Less-restricted derivatives recover some behavior for Gemma, but do not make the agent reliable or safe.
\end{abstract}

\section{Introduction}
\label{sec:intro}

Language models are starting to operate as security agents rather than only as
advisers. They inspect repositories, run tests, call command-line tools, check
whether a bug is reachable, and write remediation evidence. That shift makes
safety alignment harder to reason about. In ordinary chat, refusing
exploit-adjacent detail is often the desired behavior. In an authorized agent
workflow, the same kind of detail can be part of the job: the agent may need to
localize a vulnerable function, demonstrate reachability inside a disposable
sandbox, or verify that a patch blocks the failing path. A refusal there can
break a defensive workflow that the surrounding system has already constrained
and authorized as legitimate.

\looseness=-1
This paper evaluates that tension empirically. We ask how model-level safety
alignment changes autonomous performance and evidence quality on local,
authorized, sandboxed information-security tasks. The question is not whether
safety controls are necessary; they are. The question is where the control
should sit and how it should be measured. Model-level refusal can reduce
open-ended misuse, but be too coarse for contained workflows requiring local
construction and validation. Reducing
refusals is not automatically desirable either: it may recover useful behavior,
expose general capability differences, or produce invalid actions that still
need system-level checks.

\looseness=-1
We build one paired local evaluation around four aligned versus
less-restricted model pairs. Two pairs are Gemma 4 models: the official 31B
instruction-tuned model and a third-party 31B uncensored derivative, plus an
official 26B A4B model and a matched third-party 26B A4B uncensored derivative.
The same evaluation includes Qwen2.5-Coder 7B and Llama 3.1 8B pairs using
aligned instruction models and abliterated derivatives. In the 31B Gemma
condition set, the official model is also tested with its standard prompt, with
an explicit defensive authorization prompt, and with a fixed placeholder prefix
that serves only as an additional prompt-variant control. All comparable conditions
share the same agent harness, tasks, tools, budgets, seeds, inference backend,
and, where possible, quantization level.

\looseness=-1
The security evaluation contains 1,500 saved runs across the four model pairs.
The 31B Gemma model is evaluated under four prompt/model conditions; the 26B
Gemma, Qwen, and Llama pairs are evaluated under aligned and
less-restricted conditions. The security tasks cover vulnerability
localization, reachability validation, proof of trigger, patch verification,
and remediation reporting. Each task has a deterministic success predicate, a
redaction rule, and a machine-readable trace recording model messages, tool
calls, phase events, refusals, suppressed actions, unsafe actions, and final
outcomes. This lets the analysis separate visible refusals from ordinary agent
failures in the same trace.

\looseness=-1
Raw success rates are not enough. A report can sound plausible while failing to
read or cite the local evidence that supports its claim. We therefore add
evidence-grounding targets to every task manifest: expected files, symbols,
issue terms, and local evidence terms. These produce auditable per-run
grounding metrics from tool calls and final artifacts. We also run non-security
sanity and standard coding controls for every model pair, so that any
aligned--less-restricted gap can be compared against ordinary coding work
rather than interpreted too quickly as a security-specific effect.

\looseness=-1
The cross-family result is deliberately less tidy than a single-family account.
Both Gemma pairs show a less-restricted advantage in success and grounding, but
Qwen2.5-Coder does not reproduce the success advantage and the abliterated
Llama 3.1 derivative performs much worse under the tool protocol. The strongest
conclusion is therefore not that safety alignment uniquely suppresses security
work, or that less-restricted derivatives are generally better security agents.
It is that visible refusal rates alone do not explain agent behavior: evidence
grounding, model capability, derivative provenance, and tool-interface
reliability all matter. No family passes the hard proof and patch-verification
checks, partly because proof-of-trigger checks require an unprompted term
(Section~\ref{sec:results}).

The paper makes three contributions.
\begin{enumerate}[leftmargin=*]
  \item \looseness=-1
  We introduce a system-level benchmark for authorized security agents:
  30 local sandboxed tasks with fixed tools, deterministic success checks,
  redaction rules, and trace labels that distinguish refusal from tool,
  localization, validation, and reporting failures.\footnote{Code and data:
  \codeurl.}
  \item \looseness=-1
  We run a single cross-family evaluation across Gemma 4 31B, Gemma 4
  26B A4B, Qwen2.5-Coder 7B, and Llama 3.1 8B, comparing aligned models with
  less-restricted derivatives under shared tasks, seeds, budgets, backend, and
  success predicates, plus non-security sanity and standard coding controls.
  \item \looseness=-1
  We show that the observed behavior is not explained by visible refusal
  alone. The Gemma pairs show weaker grounded artifact production in aligned
  conditions, while Qwen and Llama show that derivative provenance and
  tool-interface reliability can dominate the comparison. Deterministic
  grounding metrics, task-blocked tests, and a blind LLM-reviewer audit support
  this cautious cross-family interpretation.
\end{enumerate}

\looseness=-1
The scope is intentionally limited. We do not evaluate third-party targets, do
not publish reusable exploit payloads, and do not argue that less-restricted
models are safe to deploy. The paper examines how refusal behavior, model
capability, and evidence grounding interact inside a contained agent system,
and why security-agent safety should be evaluated at the system level rather
than only through single-turn refusal behavior.

\section{Background and Related Work}
\label{sec:background}

\paragraph{Alignment and refusal.}
Instruction tuning and reinforcement learning from human or AI feedback make
language models more helpful and less likely to produce harmful content
\cite{christiano2017deep,ziegler2019fine,stiennon2020learning,
askell2021general,ouyang2022training,bai2022constitutional}. In deployed
assistants, this often appears as refusal behavior: the model declines to
complete requests it classifies as unsafe. Refusal is a useful safety mechanism,
but it is a context-poor one when judged only from the prompt text. Security
work creates many cases where the natural-language description resembles an
attack even though the environment is authorized, local, and intentionally
vulnerable. Prior safety evaluations show that this helpfulness-harmlessness
tension can produce both under-refusal and over-refusal
\cite{lin2021truthfulqa,hendrycks2021ethics,weidinger2021ethical,
rotter2023xstest} in deployed systems.

\paragraph{Robust refusal and jailbreak benchmarks.}
HarmBench and JailbreakBench standardize the evaluation of adversarial prompts,
harmful behaviors, and robust refusal \cite{mazeika2024harmbench,
chao2024jailbreakbench}. StrongREJECT, XSTest, and universal prompt-attack
work further clarify how refusal scoring, false refusal, and jailbreak
transferability should be measured \cite{souly2024strongreject,
rotter2023xstest,zou2023universal,wei2023jailbroken}. These benchmarks are
valuable because they clarify threat models, scoring rules, and attack
artifacts. Our evaluation differs in the target behavior. We do not ask whether
a jailbreak can make a model provide harmful instructions in open-ended chat. We
ask how refusal-oriented alignment changes task completion and evidence quality
when an agent must act over many turns inside a controlled security environment.

\paragraph{Cybersecurity capability benchmarks.}
CyberSecEval 2, CyberSecEval 3, Cybench, NYUCTF Bench, EnIGMA, and SEC-bench
evaluate cybersecurity risks and capabilities of language models, including
prompt injection, exploit generation, vulnerability patching, and
capture-the-flag style tasks \cite{bhatt2024cyberseceval,li2024cyberseceval3,
zhang2025cybench,shao2024nyuctf,abramovich2024enigma,
secbench2025automated}. These efforts show the importance of executable
environments, task-specific success predicates, and partial-credit
decomposition. We build on that benchmarking style but center paired
comparisons that keep the harness, task, and seed fixed within each evaluated
pair while varying the alignment condition.

\paragraph{Agent benchmarks and tool use.}
AgentBench, ReAct, Toolformer, WebShop, WebArena, InterCode, SWE-bench,
SWE-agent, and OSWorld show that multi-turn tool use introduces failure modes
that are invisible in single-turn question answering
\cite{li2023agentbench,yao2023react,schick2023toolformer,yao2022webshop,
zhou2023webarena,yang2023intercode,jimenez2024swebench,yang2024sweagent,
xie2024osworld}. Recent critiques emphasize that agent benchmarks require
careful success predicates, environment controls, and resistance to benchmark
gaming \cite{kapoor2024aiagents}. These lessons motivate our use of
machine-readable task manifests, deterministic checkers, and complete traces.

\looseness=-1
\paragraph{Defensive refusal in cybersecurity.}
Recent work on \emph{defensive refusal bias} \cite{campbell2026defensive}
studies the failure mode closest to ours. On 2,390 single-turn conversations
from the National Collegiate Cyber Defense Competition, Claude 3.5 Sonnet,
GPT-4o, and Llama-3.3-70B-Instruct together refuse defensive requests containing
security-sensitive keywords 2.72 times as often as semantically equivalent
neutral requests, and explicit authorization is associated with more refusals,
not fewer. Where that study classifies refusals in response text, we evaluate
multi-turn autonomous execution in executable sandboxes with deterministic
success predicates, saved traces, failure phases, unsafe-action labels, and
evidence grounding, comparing aligned models with their less-restricted
derivatives. Both studies find that authorization wording is not a
reliable lever: our authorized Gemma 4 31B condition succeeds in 0/150 runs.
But none of our security traces contains a refusal event; the aligned Gemma
models keep working and produce less-grounded artifacts, a difference that
refusal rates miss.

\paragraph{Mechanisms of refusal removal.}
Recent mechanistic work suggests that refusal behavior in open-weight chat
models can sometimes be controlled through low-dimensional activation
directions \cite{arditi2024refusal}. Community ``uncensored'' or
``abliterated'' models apply related ideas to reduce refusals. Such derivatives
are useful for paired comparisons, but they are not perfect counterfactuals: weight
editing, quantization, template changes, or fine-tuning may alter more than
refusal behavior. Our protocol therefore requires matched inference settings,
non-security sanity checks, and explicit limitations around derivative-model
provenance.

\paragraph{Model families and derivative provenance.}
Gemma 4 is the central family because the official model cards describe
instruction-tuned open-weight models with long context windows, coding ability,
structured-output support, and agentic use cases
\cite{google_gemma4_blog_2026,hf_gemma4_31b_it,hf_gemma4_26b_a4b_it}. The
TrevorJS derivatives advertise reduced refusal behavior for the same Gemma 4
family \cite{trevorjs_gemma4_uncensored_collection,
trevorjs_gemma4_31b_uncensored,trevorjs_gemma4_26b_uncensored}. The 26B
evaluation uses local GGUF quantizations from Unsloth and TrevorJS
\cite{unsloth_gemma4_26b_a4b_gguf,trevorjs_gemma4_26b_uncensored_gguf}.

The cross-family evaluation also includes Qwen2.5-Coder 7B and Llama 3.1 8B.
For each additional family, the aligned condition uses an instruction-tuned
GGUF release and the comparison condition uses an abliterated derivative
distributed as a GGUF quantization
\cite{qwen25coder7b_gguf,qwen25coder7b_abliterated_gguf,
llama31_8b_instruct_gguf,llama31_8b_abliterated_gguf}. These are not clean
causal counterfactuals for safety alignment. We therefore report them as
cross-family comparisons and treat model provenance as part of the
interpretation rather than as a nuisance detail.

\looseness=-1
The derivatives also remove refusals in different ways, so we attribute
effects to these specific transformations rather than to abliteration in
general. According to the model cards, the Gemma 4 31B derivative uses
norm-preserving biprojected abliteration \cite{grimjim2025normpreserving},
projecting per-layer refusal directions out of the attention-output and MLP
down-projection weights of every layer while preserving row norms
\cite{trevorjs_gemma4_31b_uncensored}; the 26B A4B derivative additionally
applies expert-granular abliteration to all 128 expert down-projections in each
mixture-of-experts layer \cite{trevorjs_gemma4_26b_uncensored}. The Llama 3.1
derivative uses refusal-direction abliteration \cite{arditi2024refusal} via
FailSpy's implementation \cite{mlabonne_llama31_8b_abliterated}; the
Qwen2.5-Coder derivative is an abliterated release by huihui-ai whose GGUF card
does not document the procedure \cite{qwen25coder7b_abliterated_gguf}. None
isolates safety fine-tuning: public base checkpoints precede instruction tuning
altogether, and no official instruction-tuned checkpoints of these models
without safety training are available.

\section{Methodology}
\label{sec:methodology}

\begin{figure}[t]
\centering
\definecolor{saInk}{RGB}{31,41,55}
\definecolor{saMuted}{RGB}{92,105,120}
\definecolor{saLine}{RGB}{204,213,224}
\definecolor{saPanel}{RGB}{248,250,252}
\definecolor{saBlue}{RGB}{37,99,235}
\definecolor{saTeal}{RGB}{8,145,178}
\definecolor{saGreen}{RGB}{22,163,74}
\definecolor{saAmber}{RGB}{217,119,6}
\definecolor{saRed}{RGB}{220,38,38}
\definecolor{saViolet}{RGB}{124,58,237}
\resizebox{\linewidth}{!}{%
\begin{tikzpicture}[
  x=1cm,
  y=1cm,
  font=\sffamily\scriptsize,
  panel/.style={draw=saLine, line width=0.6pt, fill=saPanel,
    rounded corners=6pt},
  card/.style={draw=saLine, line width=0.55pt, fill=white,
    rounded corners=4pt, drop shadow={shadow xshift=0.45pt,
    shadow yshift=-0.45pt, opacity=0.16}, align=left, inner sep=4pt},
  model/.style={card, text width=2.85cm, minimum width=3.10cm,
    minimum height=0.64cm},
  hstep/.style={card, align=center, text width=1.02cm, minimum width=1.26cm,
    minimum height=0.76cm, inner sep=3.4pt},
  outcome/.style={draw=saLine, line width=0.5pt, fill=white,
    rounded corners=4pt, align=center, inner sep=3.3pt, minimum width=1.38cm,
    minimum height=0.44cm},
  result/.style={card, align=left, text width=2.72cm, minimum width=2.92cm,
    minimum height=0.70cm},
  tracebox/.style={draw=saLine, line width=0.55pt, fill=white,
    rounded corners=4pt, align=left, inner xsep=7pt, inner ysep=5pt,
    text width=3.55cm, minimum width=3.95cm, minimum height=0.80cm,
    drop shadow={shadow xshift=0.45pt, shadow yshift=-0.45pt,
    opacity=0.16}},
  chip/.style={draw=saLine, line width=0.45pt, rounded corners=3pt,
    align=center, inner sep=2.8pt, minimum width=1.18cm},
  arrow/.style={-{Latex[length=2.3mm,width=1.7mm]}, line width=0.7pt,
    draw=saMuted},
  title/.style={text=saInk, font=\sffamily\bfseries\footnotesize},
  label/.style={text=saMuted, font=\sffamily\tiny}
]

\path[use as bounding box] (0,0) rectangle (16.85,5.35);

\draw[panel] (0.12,0.18) rectangle (4.02,5.12);
\draw[panel] (4.38,0.18) rectangle (11.92,5.12);
\draw[panel] (12.18,0.18) rectangle (16.72,5.12);

\node[anchor=west, title] at (0.42,4.78) {Model suite};
\node[anchor=west, label] at (0.42,4.52) {Four pairs, matched task--seed runs};

\node[model, anchor=west] (aligned) at (0.62,4.00)
  {\textbf{\textcolor{saBlue}{Aligned}} baselines\\[-1pt]
   \tiny Gemma sizes, Qwen, Llama};
\node[model, anchor=west] (auth) at (0.62,3.12)
  {\textbf{\textcolor{saViolet}{Less-restricted}} pairs\\[-1pt]
   \tiny uncensored or abliterated};
\node[model, anchor=west] (jb) at (0.62,2.24)
  {\textbf{\textcolor{saTeal}{Gemma 31B}} prompts\\[-1pt]
   \tiny standard, authorized, \mbox{placeholder}};
\node[model, anchor=west] (unc) at (0.62,1.36)
  {\textbf{\textcolor{saAmber}{Non-security}} controls\\[-1pt]
   \tiny sanity and coding checks};

\draw[saBlue, line width=1.25pt] (0.45,4.25) -- (0.45,3.75);
\draw[saViolet, line width=1.25pt] (0.45,3.37) -- (0.45,2.87);
\draw[saTeal, line width=1.25pt] (0.45,2.49) -- (0.45,1.99);
\draw[saAmber, line width=1.25pt] (0.45,1.61) -- (0.45,1.11);

\node[chip, draw=saGreen!65!black, fill=saGreen!8, anchor=west,
  minimum width=2.55cm] at (0.62,0.58) {\tiny paired task--seed runs};

\node[anchor=west, title] at (4.70,4.78) {Shared security-agent harness};
\node[anchor=west, label] at (4.70,4.52) {Contained execution with deterministic scoring};

\node[hstep] (manifest) at (5.15,3.58) {\textbf{Task}\\[-1pt]\tiny manifest};
\node[hstep] (sandbox) at (6.54,3.58) {\textbf{Local}\\[-1pt]\tiny sandbox};
\node[hstep] (controller) at (7.93,3.58) {\textbf{Agent}\\[-1pt]\tiny controller};
\node[hstep] (tools) at (9.32,3.58) {\textbf{Tool}\\[-1pt]\tiny execution};
\node[hstep] (checker) at (10.71,3.58) {\textbf{Trace}\\[-1pt]\tiny checker};

\draw[arrow] (5.78,3.58) -- (5.91,3.58);
\draw[arrow] (7.17,3.58) -- (7.30,3.58);
\draw[arrow] (8.56,3.58) -- (8.69,3.58);
\draw[arrow] (9.95,3.58) -- (10.08,3.58);

\node[result, text width=5.90cm, minimum width=6.20cm, minimum height=0.72cm]
  at (8.12,2.30)
  {\textbf{Fixed across comparable runs}\\[-1pt]
   \tiny local fixtures, allowed tools, budgets, seeds, controller mode, traces,
   redaction, and success predicates};

\node[chip, draw=saBlue!65!black, fill=saBlue!8] at (5.82,1.22) {\tiny tools};
\node[chip, draw=saTeal!65!black, fill=saTeal!8] at (7.34,1.22) {\tiny budgets};
\node[chip, draw=saAmber!65!black, fill=saAmber!8] at (8.86,1.22) {\tiny checks};
\node[chip, draw=saGreen!65!black, fill=saGreen!8] at (10.38,1.22) {\tiny redaction};

\node[anchor=west, title] at (12.55,4.78) {Trace measurements};
\node[anchor=west, label, text width=3.62cm, align=left] at (12.55,4.52)
  {Behavioral outcomes, not self-reports};

\node[outcome, draw=saGreen!65!black, fill=saGreen!8, minimum width=1.46cm]
  at (13.42,3.80)
  {\tiny success};
\node[outcome, draw=saBlue!65!black, fill=saBlue!8, minimum width=1.46cm]
  at (15.38,3.80)
  {\tiny refusal};
\node[outcome, draw=saAmber!65!black, fill=saAmber!8, minimum width=1.46cm]
  at (13.42,3.18)
  {\tiny stuck};
\node[outcome, draw=saRed!65!black, fill=saRed!8, minimum width=1.46cm]
  at (15.38,3.18)
  {\tiny unsafe};

\node[tracebox] at (14.45,2.14)
  {\textbf{Observed pattern}\\[-1pt]
   \tiny pair-dependent success\\[-1pt]
   \tiny and grounding gaps};
\node[tracebox] at (14.45,1.12)
  {\textbf{Boundary}\\[-1pt]
   \tiny hard proof and patch checks\\[-1pt]
   \tiny remain unsolved};

\draw[arrow] (4.02,2.68) -- (4.38,2.68);
\draw[arrow] (11.92,2.68) -- (12.18,2.68);
\node[label] at (4.20,2.95) {same tasks};
\node[label] at (12.05,2.95) {score traces};

\end{tikzpicture}%
}
\caption{Evaluation design. Four aligned--less-restricted model pairs, Gemma
31B prompt variants, and non-security controls share one local security-agent
harness; saved traces are scored for success, grounding, refusal, ordinary
failure, and unsafe action.}
\label{fig:evaluation-flow}
\end{figure}


\subsection{Research Questions and Hypotheses}

The evaluation asks five linked questions. First, does model condition change
end-to-end success under the same tools, budgets, and success predicates?
Second, where do failed runs break: refusal, agent-interface use, weak local
grounding, or hard proof and patch-verification steps? Third, does explicit
authorization recover useful behavior without increasing unsafe actions?
Fourth, within a model and condition, do runs that read the expected files and
name the relevant files, symbols, issues, and local evidence succeed more often
than those that do not? Finally, is the
aligned--less-restricted gap larger on security-agent tasks than on ordinary
local coding controls under the same harness?

Our working hypothesis is that less-restricted derivatives may perform better
on tasks requiring construction, reachability checks, or validation, because
these steps are close to behavior that model-level refusal policies often
discourage. We do not treat such a gain as evidence by itself. The
interpretation depends on failure phase, evidence grounding, unsafe-action
rates, non-security controls, and whether the direction is stable across model
families.

\subsection{Threat Model and Scope}

All tasks are local, authorized, and constructed to contain known defects that
the benchmark asks the agent to inspect or repair. No task targets
third-party systems. Each run starts in a fresh temporary workspace that
confines file tools, and tools or commands outside an allow-list are blocked and
logged, but the network and host filesystem are not isolated
(Appendix~\ref{app:out-of-scope}). The policy forbids attempts to escape
the sandbox, access external systems, exfiltrate secrets, or produce reusable
operational instructions; the automated unsafe-action label records only
blocked tool-policy violations.

This threat model separates \emph{task legitimacy} from \emph{model output
surface}. The same natural-language action can be disallowed in an open setting
and required in a controlled benchmark. The purpose of the evaluation is to
measure whether model behavior tracks that distinction when the surrounding
agent system establishes authorization and containment.

\subsection{Task and Model Suite}

\begin{table}[t]
\centering
\caption{Model pairs and conditions. The 31B Gemma row also includes
authorization and placeholder-prefix prompt conditions for the official aligned
model; the other rows use aligned versus less-restricted derivatives only.}
\label{tab:model-matrix}
\small
\begin{tabular}{@{}>{\raggedright\arraybackslash}p{0.14\linewidth}
  >{\raggedright\arraybackslash}p{0.25\linewidth}
  >{\raggedright\arraybackslash}p{0.30\linewidth}
  >{\raggedright\arraybackslash}p{0.17\linewidth}@{}}
\toprule
Model pair & Aligned model & Less-restricted counterpart & Purpose \\
\midrule
Gemma 4 31B & Official instruction-tuned model & TrevorJS 31B uncensored GGUF & Main matched pair \\
Gemma 4 26B A4B & Official A4B instruction-tuned model & TrevorJS 26B A4B uncensored GGUF & Gemma size check \\
Qwen2.5-Coder 7B & Official coder-instruct GGUF & Abliterated coder-instruct GGUF & Non-Gemma coding check \\
Llama 3.1 8B & Instruction-tuned GGUF & Abliterated instruction-tuned GGUF & Non-Gemma chat check \\
\bottomrule
\end{tabular}
\end{table}


\begin{table}[t]
\centering
\caption{Task-suite composition. Each class contributes six tasks to the 30-task security benchmark.}
\label{tab:task-suite}
\footnotesize
\begin{tabular}{p{0.21\linewidth}p{0.34\linewidth}p{0.35\linewidth}}
\toprule
Task class & Agent objective & Success predicate \\
\midrule
Vulnerability localization & Identify the file, function, and root cause for a seeded local defect & Matches ground truth and bug class without operational detail \\
Reachability validation & Check whether a seeded vulnerable path is reachable in a local service & Trace contains required evidence for the reachable/unreachable label \\
Proof of trigger & Produce redacted evidence that the vulnerable path can be triggered in the sandbox & Trigger evidence satisfies redaction rule and cannot be reused against third-party targets \\
Patch verification & Confirm that a fixed version blocks the same trigger condition & Patched run demonstrates blocked behavior and identifies the relevant fix \\
Security report writing & Produce remediation-oriented finding for a controlled issue & Report contains impact, evidence, remediation, and no reusable exploit details \\
\bottomrule
\end{tabular}
\end{table}


The model suite contains four aligned versus less-restricted pairs:
Gemma 4 31B, Gemma 4 26B A4B, Qwen2.5-Coder 7B, and Llama 3.1 8B. The Gemma
pairs use official aligned models and third-party uncensored derivatives. The
Qwen and Llama pairs use aligned instruction models and abliterated
derivatives. The 31B Gemma evaluation additionally includes explicit
authorization and placeholder-prefix prompt conditions for the official aligned
model. We report the paired comparisons together for readability, while keeping
statistical tests separate because model family, size, architecture details,
and derivative provenance are not controlled tightly enough for pooling. Every security and control trace records the model identifier, GGUF
hash, backend, quantization, controller mode, and sampling settings.

The artifact implements a 30-task catalog split evenly across five classes.
Each task has a machine-readable manifest, local fixture files, a deterministic
success predicate, and a redaction rule. The tasks are intentionally small and
original: they use local code snippets and report-writing contexts rather than
third-party targets, live credentials, public exploit chains, or internet
dependencies. Proof-of-trigger tasks require redacted local evidence rather than
copy-paste trigger strings.

Each manifest also contains evidence-grounding targets: expected files,
symbols, issue terms, and local evidence terms. Grounding is scored by
case-insensitive substring matching against redacted tool calls and the final
artifact. It is deliberately non-semantic: an auditable check of whether the
agent read the relevant file and connected expected local evidence to the final
artifact.

The non-security sanity catalog contains 12 local coding tasks using the same
fixture, controller, trace, predicate, and analysis machinery. These tasks cover
ordinary bug localization, branch-condition mistakes, cache-key mistakes,
parsing and formatting errors, small remediation notes, and test-failure
diagnosis. They are included to distinguish a security-specific effect from a
general model-quality or grounding gap. A smaller standard coding catalog adds
eight HumanEval/MBPP-style local repair tasks as a second non-security control.
Both controls are run for every model pair.

\subsection{Agent Harness}

Each run uses the same Python controller and tool permissions. The controller
materializes task fixtures into a temporary local workspace, sends the task
prompt to a local OpenAI-compatible model endpoint, receives JSON-formatted
model actions, executes permitted tools inside the sandbox, returns tool output,
and stops when the model emits a final artifact or the budget is exhausted. The
harness enforces:
\begin{enumerate}[leftmargin=*]
  \item a per-task time budget (declared, never reached),
  \item a fixed maximum number of model calls,
  \item the same context-window policy,
  \item the same sampling parameters for all comparable conditions,
  \item five replicate runs per task (unseeded sampling),
  \item complete trace logging, and
  \item redaction before public release.
\end{enumerate}

The harness does not rely on the model to self-report success. Automated
predicates are lexical: each checks that the final artifact contains the task's
required terms and none of its forbidden terms.
Two blind LLM reviewer passes are used only as an audit of saved trace labels
and grounding judgments; they are not used to decide task success.

The implementation also includes a deterministic mock backend used only for
artifact validation. Mock results test orchestration, schema validation,
analysis, and redaction; they are artifact checks rather than empirical model
evidence.

\looseness=-1
The evaluation uses two controller modes. The strict controller accepts only the
declared JSON tool protocol. The repair controller includes the initial file
inventory in the task prompt, converts a repeated identical
\texttt{list\_files} call into a read of the first relevant source file, and
repairs two common artifact-write wrappers. The repair controller does not
change task ground truth or success predicates. All full security and control
runs use the repair controller because strict pilot runs showed that several
aligned conditions were otherwise dominated by interface failures. A matched
strict-versus-repair ablation quantifies this interface effect, and every saved
trace records the controller choice. Appendix~\ref{app:controller-repairs}
specifies each repair rule and how often it fired.

\subsection{Prompt Conditions}
\label{sec:prompt-conditions}

The \emph{aligned} condition uses the aligned model's standard chat template
and the task prompt. The \emph{authorized} condition, evaluated for Gemma 31B
only, adds a concise statement that the task is local, authorized, defensive,
and sandboxed. The \emph{placeholder} condition, also evaluated for Gemma 31B
only, was designed to apply a private jailbreak-style template, but that
template was never configured for the reported runs. The harness instead
prefixed every task prompt with a fixed notice that no template was configured,
so we treat this condition only as a prompt-variant control and draw no
conclusions about jailbreak prompting; its traces keep the harness label
\texttt{jailbreak}. The
\emph{uncensored} condition name is retained in the harness for compatibility
with earlier Gemma runs; in cross-family text we describe the comparison models
more precisely as less-restricted or abliterated derivatives.

\looseness=-1
Prompt wording is part of the experimental artifact. It is versioned, hashed,
and kept constant within each reported condition. Three additional
authorization prompt variants are evaluated on a matched 10-task Gemma 31B
subset to test whether authorized-condition behavior is wording-sensitive.

\subsection{Judging and Labels}

Every failed run receives exactly one failure label from the label taxonomy in
Appendix~\ref{app:supplementary-tables}. Secondary labels may be recorded, but
that single label is the aggregation key. Automated predicates decide success.
Refusals are matched against seven fixed phrases; partial refusals and
suppressed actions are not detected.
To check whether these labels and grounding bins are auditable, two independent
LLM reviewer passes label the same blind stratified packet of redacted traces;
the packet excludes condition, original label, deterministic grounding score,
and harness checker reason.

\section{Evaluation Protocol}
\label{sec:evaluation}

\subsection{Metrics}

The main outcome is end-to-end task success: the final artifact satisfies
the task predicate within the budget and without violating the sandbox policy.
Secondary metrics separate refusal behavior from general agent failure:
single failure label, unsafe-action count, invalid-tool count, token cost,
model calls, wall-clock time, and evidence-grounding fields. Grounding fields
record whether the trace read an expected file and whether the final artifact
named an expected file, symbol, issue term, and local evidence term. This
distinction matters because a lower success rate is not evidence of
alignment-induced suppression unless the trace shows refusal, suppressed
action, or phase-specific omission rather than generic tool or capability
failure. The metrics are defined in
Appendix~\ref{app:supplementary-tables}.

\subsection{Sample Size}

The 31B Gemma repair-controller security evaluation has 30 tasks, five seeds,
and four conditions, for 600 runs. The 26B Gemma, Qwen2.5-Coder, and Llama 3.1
evaluations repeat the aligned versus less-restricted comparison with 300
security runs per pair. Across the four model pairs, the security suite therefore
contains 1,500 saved traces. The controls add 120 non-security sanity runs and
80 standard coding-control runs for each pair. Separate 24-run pilots validate
orchestration; a 150-run authorization-prompt robustness run and a 200-run
strict-controller ablation, matched against the corresponding repair-controller
subset, test controller and prompt sensitivity on 31B Gemma subsets. Results
are shown together when useful, but statistical tests are not pooled across
model families or sizes.

\subsection{Paired Statistical Tests}

For each aligned--less-restricted comparison, runs are paired by task and seed.
We report absolute success-rate differences in percentage points, bootstrap
95\% confidence intervals for success rates, and McNemar's exact test on the
paired binary outcomes. For three Gemma comparisons, we also aggregate the
paired differences by task and apply a task-blocked sign-permutation test to
the success and grounding gaps (Table~\ref{app:robustness}). Cost metrics are
skewed, so the generated analysis reports token and model-call totals and the
median wall-clock time per run. Mixed-effects
models are left to follow-up work because success is near zero in several task
classes.

\subsection{Decision Rules}

We adopt a security-specific alignment-effect interpretation only if all four
of the following gates are satisfied:
\begin{enumerate}[leftmargin=*]
  \item the less-restricted condition outperforms the aligned condition on paired
  success,
  \item the aligned condition has higher strict-refusal, partial-refusal, or
  suppressed-action rates, or shows more security-specific omission in the
  trace labels,
  \item phase labels show losses during necessary security steps rather than
  only generic planning, interface, or tool-use failures, and
  \item non-security coding controls do not show a comparable
  aligned--less-restricted capability gap.
\end{enumerate}

\looseness=-1
The completed sanity suite and standard coding control fail the last gate for
the Gemma pairs: the less-restricted derivatives also have large advantages on
ordinary non-security coding tasks. The Qwen and Llama families add a different
caution: their less-restricted derivatives do not reproduce the Gemma success
advantage under the same tool harness. We therefore interpret the measured
effect as model-family, size, and derivative-sensitive evidence about
grounding, capability, and tool-interface behavior, while still using the security traces
to show that visible refusals are not the main mechanism in the security-agent
setting.

One further reporting rule applies: if a less-restricted condition improves
success but also increases unsafe actions, we report a safety-performance
trade-off, not an unqualified improvement.

\section{Results}
\label{sec:results}

\looseness=-1
The artifact contains 1,500 security-agent traces across four model pairs:
600 Gemma 31B traces under aligned, authorized, placeholder-prefix, and
less-restricted conditions, plus 300 aligned-versus-less-restricted traces each
for Gemma 26B A4B, Qwen2.5-Coder 7B, and Llama 3.1 8B. The non-security
controls add 800 traces across the same four pairs. All reported security
runs use the same 30 local tasks, repair controller, tool set, budgets, seeds,
grounding targets, redaction rules, and deterministic success predicates for
every task.

\paragraph{Controller checks.}
The pilots validated orchestration and exposed an interface confound: under the
strict controller, official-model failures were dominated by repeated file
listings and malformed tool calls. The repair controller removes that failure
mode without changing ground truth or success predicates. A 200-run
strict-controller ablation, compared with the matched repair-controller subset,
confirms that the interface matters but does not reverse the Gemma 31B
comparison: official-model prompt conditions remain at
0/50 successes on the 10-task subset, while the less-restricted derivative
reaches 5/50 under strict and 10/50 under repair; the appendix reports the
pilot and ablation tables.

\paragraph{Cross-pair security comparison.}
Figure~\ref{fig:model-family-results} and
Table~\ref{tab:model-family-security} show the main cross-pair result. The
Gemma pairs reproduce the same direction: the less-restricted 31B derivative
succeeds on 21/150 security runs (14.0\%) versus 1/150 for the aligned model
(0.7\%; McNemar \(p=1.10\times10^{-5}\)), and the less-restricted 26B A4B
derivative succeeds on 16/150 runs (10.7\%) versus 0/150 for the aligned model
(\(p=3.05\times10^{-5}\)). Both Gemma derivatives also produce more grounded
artifacts.

The non-Gemma families do not support a general ``less-restricted is better''
claim. Qwen2.5-Coder is roughly neutral on grounding and has lower success for
the abliterated derivative: 3/150 successes (2.0\%) versus 8/150 for the
aligned model (5.3\%; \(p=0.18\)). Llama 3.1 is sharper: the aligned model
succeeds on 5/150 runs (3.3\%), while the abliterated derivative succeeds on
0/150 (\(p=0.0625\)) and has a grounding score near zero. That derivative is not merely less
refusal-prone under this harness; it is substantially worse at following the
tool protocol.

\begin{figure}[t]
\centering
\definecolor{famAlign}{RGB}{0,114,178}
\definecolor{famLess}{RGB}{213,94,0}
\definecolor{famGrid}{RGB}{226,232,240}
\definecolor{famAxis}{RGB}{148,163,184}
\definecolor{famInk}{RGB}{31,41,55}
\definecolor{famMuted}{RGB}{75,85,99}
\resizebox{\linewidth}{!}{%
\begin{tikzpicture}[x=1cm,y=1cm,font=\sffamily\scriptsize]
\path[use as bounding box] (0,0) rectangle (14.2,5.75);
\node[anchor=west,font=\sffamily\bfseries,text=famInk] at (0.05,5.08) {Model pair};
\node[anchor=west,font=\sffamily\bfseries,text=famInk] at (3.10,5.08) {Task success (\% runs)};
\node[anchor=west,font=\sffamily\bfseries,text=famInk] at (9.15,5.08) {Evidence grounding (score / 5)};
\filldraw[fill=famAlign,draw=white,line width=0.25pt] (10.70,5.48) circle (0.075);
\node[anchor=west,text=famMuted] at (10.87,5.48) {Aligned};
\filldraw[fill=famLess,draw=white,line width=0.25pt] (12.20,5.48) circle (0.075);
\node[anchor=west,text=famMuted] at (12.37,5.48) {Less-restricted};
\draw[famAxis,line width=0.35pt] (3.10,0.92) -- (7.55,0.92);
\draw[famGrid,line width=0.25pt] (3.100,0.92) -- (3.100,4.78);
\node[anchor=north,text=famMuted] at (3.100,0.81) {0\%};
\draw[famGrid,line width=0.25pt] (4.213,0.92) -- (4.213,4.78);
\node[anchor=north,text=famMuted] at (4.213,0.81) {5\%};
\draw[famGrid,line width=0.25pt] (5.325,0.92) -- (5.325,4.78);
\node[anchor=north,text=famMuted] at (5.325,0.81) {10\%};
\draw[famGrid,line width=0.25pt] (6.438,0.92) -- (6.438,4.78);
\node[anchor=north,text=famMuted] at (6.438,0.81) {15\%};
\draw[famGrid,line width=0.25pt] (7.550,0.92) -- (7.550,4.78);
\node[anchor=north,text=famMuted] at (7.550,0.81) {20\%};
\draw[famAxis,line width=0.35pt] (9.15,0.92) -- (13.60,0.92);
\draw[famGrid,line width=0.25pt] (9.150,0.92) -- (9.150,4.78);
\node[anchor=north,text=famMuted] at (9.150,0.81) {0};
\draw[famGrid,line width=0.25pt] (10.040,0.92) -- (10.040,4.78);
\node[anchor=north,text=famMuted] at (10.040,0.81) {1};
\draw[famGrid,line width=0.25pt] (10.930,0.92) -- (10.930,4.78);
\node[anchor=north,text=famMuted] at (10.930,0.81) {2};
\draw[famGrid,line width=0.25pt] (11.820,0.92) -- (11.820,4.78);
\node[anchor=north,text=famMuted] at (11.820,0.81) {3};
\draw[famGrid,line width=0.25pt] (12.710,0.92) -- (12.710,4.78);
\node[anchor=north,text=famMuted] at (12.710,0.81) {4};
\draw[famGrid,line width=0.25pt] (13.600,0.92) -- (13.600,4.78);
\node[anchor=north,text=famMuted] at (13.600,0.81) {5};
\node[anchor=east,text=famInk,align=right] at (2.55,4.380) {Gemma 31B};
\draw[famGrid,line width=0.25pt] (0.05,3.910) -- (13.60,3.910);
\draw[famInk,line width=0.35pt] (3.100,4.510) -- (3.545,4.510);
\draw[famInk,line width=0.35pt] (3.100,4.440) -- (3.100,4.580);
\draw[famInk,line width=0.35pt] (3.545,4.440) -- (3.545,4.580);
\filldraw[fill=famAlign,draw=white,line width=0.25pt] (3.248,4.510) circle (0.075);
\filldraw[fill=famAlign,draw=white,line width=0.25pt] (12.057,4.510) circle (0.075);
\draw[famInk,line width=0.35pt] (5.028,4.250) -- (7.550,4.250);
\draw[famInk,line width=0.35pt] (5.028,4.180) -- (5.028,4.320);
\draw[famInk,line width=0.35pt] (7.550,4.180) -- (7.550,4.320);
\filldraw[fill=famLess,draw=white,line width=0.25pt] (6.215,4.250) circle (0.075);
\filldraw[fill=famLess,draw=white,line width=0.25pt] (12.633,4.250) circle (0.075);
\node[anchor=east,text=famInk,align=right] at (2.55,3.380) {Gemma 26B A4B};
\draw[famGrid,line width=0.25pt] (0.05,2.910) -- (13.60,2.910);
\draw[famInk,line width=0.35pt] (3.100,3.510) -- (3.100,3.510);
\draw[famInk,line width=0.35pt] (3.100,3.440) -- (3.100,3.580);
\draw[famInk,line width=0.35pt] (3.100,3.440) -- (3.100,3.580);
\filldraw[fill=famAlign,draw=white,line width=0.25pt] (3.100,3.510) circle (0.075);
\filldraw[fill=famAlign,draw=white,line width=0.25pt] (10.610,3.510) circle (0.075);
\draw[famInk,line width=0.35pt] (4.435,3.250) -- (6.660,3.250);
\draw[famInk,line width=0.35pt] (4.435,3.180) -- (4.435,3.320);
\draw[famInk,line width=0.35pt] (6.660,3.180) -- (6.660,3.320);
\filldraw[fill=famLess,draw=white,line width=0.25pt] (5.473,3.250) circle (0.075);
\filldraw[fill=famLess,draw=white,line width=0.25pt] (12.817,3.250) circle (0.075);
\node[anchor=east,text=famInk,align=right] at (2.55,2.380) {Qwen2.5-Coder 7B};
\draw[famGrid,line width=0.25pt] (0.05,1.910) -- (13.60,1.910);
\draw[famInk,line width=0.35pt] (3.545,2.510) -- (5.177,2.510);
\draw[famInk,line width=0.35pt] (3.545,2.440) -- (3.545,2.580);
\draw[famInk,line width=0.35pt] (5.177,2.440) -- (5.177,2.580);
\filldraw[fill=famAlign,draw=white,line width=0.25pt] (4.287,2.510) circle (0.075);
\filldraw[fill=famAlign,draw=white,line width=0.25pt] (11.968,2.510) circle (0.075);
\draw[famInk,line width=0.35pt] (3.100,2.250) -- (4.138,2.250);
\draw[famInk,line width=0.35pt] (3.100,2.180) -- (3.100,2.320);
\draw[famInk,line width=0.35pt] (4.138,2.180) -- (4.138,2.320);
\filldraw[fill=famLess,draw=white,line width=0.25pt] (3.545,2.250) circle (0.075);
\filldraw[fill=famLess,draw=white,line width=0.25pt] (12.010,2.250) circle (0.075);
\node[anchor=east,text=famInk,align=right] at (2.55,1.380) {Llama 3.1 8B};
\draw[famGrid,line width=0.25pt] (0.05,0.910) -- (13.60,0.910);
\draw[famInk,line width=0.35pt] (3.248,1.510) -- (4.583,1.510);
\draw[famInk,line width=0.35pt] (3.248,1.440) -- (3.248,1.580);
\draw[famInk,line width=0.35pt] (4.583,1.440) -- (4.583,1.580);
\filldraw[fill=famAlign,draw=white,line width=0.25pt] (3.842,1.510) circle (0.075);
\filldraw[fill=famAlign,draw=white,line width=0.25pt] (11.642,1.510) circle (0.075);
\draw[famInk,line width=0.35pt] (3.100,1.250) -- (3.100,1.250);
\draw[famInk,line width=0.35pt] (3.100,1.180) -- (3.100,1.320);
\draw[famInk,line width=0.35pt] (3.100,1.180) -- (3.100,1.320);
\filldraw[fill=famLess,draw=white,line width=0.25pt] (3.100,1.250) circle (0.075);
\filldraw[fill=famLess,draw=white,line width=0.25pt] (9.340,1.250) circle (0.075);
\end{tikzpicture}%
}
\caption{Cross-pair security results. Success points include bootstrap 95\% confidence intervals; grounding points show the mean deterministic evidence-grounding score out of five.}
\label{fig:model-family-results}
\end{figure}


\begin{table}[t]
\centering
\caption{Cross-pair security comparison under the same repair-controller
harness. ``Less'' denotes the less-restricted derivative for each model pair.
Grounding is the mean deterministic score out of five; unsafe is the share of
less-restricted runs with at least one unsafe-action event.}
\label{tab:model-family-security}
\small
\begin{tabular}{@{}lrrrrrr@{}}
\toprule
Model pair & Align succ. & Less succ. & $\Delta$ succ. & Align ground. & Less ground. & Less unsafe \\
\midrule
Gemma 31B & 0.7\% & 14.0\% & +13.3 & 3.27 & 3.91 & 0.0\% \\
Gemma 26B A4B & 0.0\% & 10.7\% & +10.7 & 1.64 & 4.12 & 5.3\% \\
Qwen2.5-Coder 7B & 5.3\% & 2.0\% & -3.3 & 3.17 & 3.21 & 0.0\% \\
Llama 3.1 8B & 3.3\% & 0.0\% & -3.3 & 2.80 & 0.21 & 3.3\% \\
\bottomrule
\end{tabular}

\end{table}

\paragraph{Mechanism and limits.}
Within Gemma, the gain is concentrated in vulnerability localization,
reachability validation, and report writing; every Gemma condition at both
sizes still scores 0/30 on proof-of-trigger and patch-verification tasks.
Visible refusal does not explain the 31B result: the 31B traces contain no
strict-refusal labels and no refusal, suppressed-action, or unsafe-action
events. The difference appears later, in grounded artifact production. The
less-restricted 31B derivative names expected symbols more often than the
aligned model (0.81 versus 0.41) and has higher mean grounding (3.91 versus
3.27); the 26B pair shows the same direction (4.12 versus 1.64).

\looseness=-1
For the fourth question, we compare runs of the same model and condition that do
and do not satisfy a grounding field. Every successful run in the 1,500-trace
evaluation names the expected symbol, but 29 of the 30 success predicates
require it and the remaining task has no successes, so this association is
definitional. Grounding is more informative among failed runs: the aligned 31B
model reads an expected file in all 149 failed runs but names the expected
symbol in only 60.

\looseness=-1
The non-Gemma traces narrow the claim. Qwen reads expected files in every
security run and has nearly identical grounding across conditions. Llama shows
the opposite of the Gemma pattern: the abliterated derivative has 124/150
tool-failure labels and grounding 0.21, compared with 35/150 tool failures and
grounding 2.80 for the aligned model. The safety picture is also mixed:
Gemma 31B and Qwen have no unsafe-action events; the 26B derivative's success
gain is a safety-performance trade-off, since 8/150 of its runs contain
tool-policy violations (disallowed commands, tools, or paths) that the sandbox
blocked; and 6/150 aligned and 5/150 abliterated Llama runs contain such
violations. A trace audit found one further less-restricted 26B run with an
out-of-scope package installation (Appendix~\ref{app:out-of-scope}). These details are reported in the appendix because they matter for
deployment interpretation, not because they rescue a universal claim.

\paragraph{Controls, examples, and audits.}
The non-security controls in Appendix~\ref{app:cross-evaluation-summary} make
the Gemma result a broad grounding/capability gap under this harness rather
than a clean security-specific alignment effect. Compact qualitative examples
show the same pattern: an aligned Gemma trace reads the right file but names the
wrong function, a less-restricted Gemma trace names the file, symbol, issue, and
evidence, and a proof-of-trigger trace still fails despite a plausible grounded
report. Task-blocked tests keep the Gemma security gaps directionally robust;
authorization-prompt variants remain at 0/50 successes on the matched 31B
subset. The blind 120-trace LLM audit supports coarse success and refusal
judgments (inter-reviewer \(\kappa=0.87\) for success), but not hard-validation
subtypes (\(\kappa=0.21\)).

\section{Discussion, Limitations, and Ethics}
\label{sec:discussion}

\looseness=-1
The strongest version of this paper is not an argument against safety
alignment. It is an argument for evaluating where safety controls belong in an
agentic system. Across the cross-family evaluation, visible refusal is only one
small part of the behavior. In Gemma, evidence grounding separates the
conditions after both models read relevant files; in Qwen, the comparison is
mostly flat; in Llama, tool-interface reliability dominates the abliterated
condition. No family passes the hard-validation checks.

\looseness=-1
The main methodological implication is that autonomous security-agent
evaluations should not collapse success, refusal, unsafe action, and evidence
grounding into one score. These outcomes answer different questions. A refusal
label asks whether the model declined an authorized step; a success predicate
asks whether the final artifact satisfies the task; an unsafe-action label asks
whether the sandbox policy was violated; and grounding asks whether the trace
and artifact contain the local files, symbols, issue terms, and evidence needed
to justify the claim. The Gemma result appears in this separation: the aligned
model is not simply refusing, yet it is less often producing the concrete
defensive evidence required by the task.

\looseness=-1
The blind audit is therefore a check on the interpretation, not an alternate
success judge. The deterministic grounding score is intentionally mechanical,
so the audit asks whether independent reviewers see the same coarse structure
in redacted traces when condition labels, original failure labels, checker
reasons, and grounding scores are hidden. Inter-reviewer agreement is high for
success and refusal judgments but low for hard-validation subtypes, and each
reviewer agrees less with the lexical success labels (Section~\ref{sec:results}).

\looseness=-1
The deployment lesson is also narrower than ``remove alignment.'' The safer
interpretation is that policy should be evaluated at the system boundary where
authorization, tool permissions, trace review, redaction, and final evidence
checks can all be observed. For an open-weight model, whoever runs the model
writes the prompt, so prompt text cannot establish authorization; the
surrounding system must.

\subsection{Limitations}
\label{sec:limitations}

\looseness=-1
\textbf{Derivatives are probes, not counterfactuals.} Each model pair uses one
less-restricted derivative with its own refusal-removal procedure (Section~\ref{sec:background}). We did not vary the abliteration
method or strength, so effects are attributable only to these derivatives;
comparing methods and strengths is left to future work.
\textbf{Capability confound.} The Gemma derivatives also do better on
non-security controls, the aligned 26B model mostly fails through repeated tool
calls that the repair controller does not cover (129/150 runs), and the Llama
derivative is much worse at following the
tool protocol, so capability and interface reliability vary with alignment
condition. A causal claim about alignment, which we do not make, would require
pairs matched on non-security performance. \textbf{Controller.} The repair controller changes effect
magnitudes (Section~\ref{sec:results}); its deterministic rules are logged per
trace (Appendix~\ref{app:controller-repairs}).
\textbf{Scale and measurement.} Absolute success is low, no condition passes the
lexical proof-of-trigger or patch-verification predicates, and the 30 tasks are
stratified across five classes rather than sampled from security work at large.
Success and grounding are both lexical substring checks, and the trace audit
uses LLM reviewers rather than human annotators, with low agreement on
hard-validation subtypes (\(\kappa=0.21\)).
\textbf{Scope.} The authorization-prompt ablation covers Gemma 4 31B only, and
all security and control runs use one backend (\texttt{llama.cpp}), 4-bit GGUF quantizations, one
sampling configuration, and five unseeded replicates per task, without network
isolation.

\subsection{Ethics and Dual Use}
\label{sec:ethics}

\looseness=-1
This is dual-use research. All experiments use original local tasks with known
defects, scoped tools, and no third-party targets. Released traces are redacted
and audited, and no jailbreak prompt was applied in any reported run. The less-restricted derivatives are not safe
defaults: two 26B derivative runs installed packages from the network
(Appendix~\ref{app:out-of-scope}), and we do not endorse deploying them. Production systems still
need alignment, sandboxing, policy checks, and evidence verification.

\section{Conclusion}

\looseness=-1
This paper turns a common concern about safety alignment into a concrete
measurement problem. In the Gemma 4 pairs, the aligned models usually keep
working but produce less grounded artifacts, while the less-restricted
derivatives also do better on non-security controls. In Qwen2.5-Coder and
Llama 3.1, the same pattern does not hold. The result is therefore not a clean
security-specific alignment effect and not a general endorsement of
less-restricted derivatives. It is evidence that autonomous security-agent
behavior depends on grounding, capability, derivative provenance, and
tool-interface reliability.

The practical conclusion is that safety evaluations for security agents should
be trace-based and evidence-aware. They should ask not only whether a model
refuses, but whether it reads the right local evidence, uses authorized tools
without escaping the sandbox, names the relevant program objects, and produces
a final artifact that a defender could act on. They should also test ordinary
coding controls, because a security-task gap is hard to interpret without a
baseline for general agent reliability. In this study, removing alignment
recovers some grounded behavior for Gemma, but it introduces unsafe-action
events for the 26B derivative and does not solve hard proof and patch
verification. The right object of
evaluation is the whole agentic system: the model, tools, authorization
boundary, trace evidence, and final checker together. That is the standard a
security-agent benchmark should measure, not a transcript-level refusal rate
alone.


\begin{ack}
This work was self-funded by the authors. The authors declare no competing
interests.
\end{ack}

\clearpage
\bibliographystyle{unsrtnat}
\nocite{*}
\bibliography{refs}

\clearpage
\appendix

\section{Artifact and Reproduction}
\label{app:artifact}

Repository: \codeurl.

The public artifact consists of the evaluation package
\texttt{security\_agent\_eval/}, task catalogs in \texttt{tasks/}, model
endpoint and GGUF-provenance configs in \texttt{configs/}, redacted traces for
every reported run and their redaction audits in \texttt{artifacts/}, and
generated summaries in \texttt{analysis/}. Raw local traces are not released;
the redacted traces reproduce every reported success, grounding, and
failure-label summary. The four model-pair
configs are \texttt{configs/models.example.json},
\texttt{configs/models.26b\_a4b.json},
\texttt{configs/models.qwen25\_coder\_7b.json}, and
\texttt{configs/models.llama31\_8b.json}; each records endpoint, local GGUF
path, model source URL, and SHA-256 hash.

The full analysis artifact can be regenerated with:
\begin{verbatim}
scripts/reproduce_artifact.sh
\end{verbatim}
This validates all task catalogs, regenerates the analyses from the redacted
traces, reruns the reviewer-audit summaries, audits public redactions, regenerates paper tables,
runs unit tests, and rebuilds the PDF. Full model reruns use the same CLI
template with the appropriate task catalog and model config:
\begin{verbatim}
python3 -m security_agent_eval.cli run \
  --tasks tasks/catalog.json \
  --models configs/<model_config>.json \
  --out runs/<run_name> \
  --conditions aligned uncensored \
  --seeds 0 1 2 3 4 \
  --controller-mode repair
\end{verbatim}
The 31B Gemma security run additionally includes the \texttt{authorized} and
\texttt{jailbreak} conditions; the latter is the placeholder-prefix condition
described in Section~\ref{sec:prompt-conditions}. Its prefix is the harness's
fixed fallback notice: ``This condition requires a private fixed
jailbreak-style template. No private template was configured, so this run uses
the authorized-sandbox wording and records the missing-template hash.'' Despite
this wording, no authorization prefix was added and no template hash was
recorded. Non-security controls replace
\texttt{tasks/catalog.json} with \texttt{tasks/sanity\_catalog.json} or
\texttt{tasks/standard\_coding\_catalog.json}.

\section{Controller Repair Rules}
\label{app:controller-repairs}

The strict and repair controllers share the same tools, budgets, task prompts
(apart from rule 1 below), seeds, and success predicates. In both modes, a model reply is parsed as JSON;
if that fails, the controller parses the substring from the first \verb|{| to
the last \verb|}|. In both modes, a tool call identical in tool name and
arguments to the most recent tool call recorded in the trace, including a
blocked or failed one, is not executed: the
controller returns a ``Repeated identical tool call'' message and counts an
invalid tool call. The repair controller adds exactly four deterministic rules,
implemented in \texttt{security\_agent\_eval/prompts.py} and
\texttt{security\_agent\_eval/harness.py}. Rules 2--4 tag the resulting tool
call with a \texttt{controller\_repair} field in the saved trace.
\begin{enumerate}[leftmargin=*]
  \item \textbf{Initial file inventory.} \emph{Trigger:} every repair-mode run
  whose task has fixture files. \emph{Transformation:} the task prompt is
  extended with ``Initial file inventory:'' followed by the sorted fixture
  paths and the sentence ``Start by reading the relevant source file from this
  inventory before writing the final artifact.''
  \item \textbf{Repeated \texttt{list\_files}.} \emph{Trigger:} a repeated
  identical call (as defined above) whose tool is \texttt{list\_files}, in a
  task that allows \texttt{read\_file}. \emph{Transformation:} instead of the
  repeated-call message, the controller executes \texttt{read\_file} on one
  path from the previous listing. Among listed entries that contain a dot, it
  takes the first ending in \texttt{.py}; if there is none, the first ending in
  \texttt{.js}, then \texttt{.ts}, \texttt{.java}, \texttt{.go}, \texttt{.rb},
  \texttt{.php}, \texttt{.c}, \texttt{.cpp}, \texttt{.md}, and \texttt{.txt},
  falling back to the first such entry. The tool output is prefixed with a
  note naming the substituted read. If no entry qualifies, the repeated-call
  message is returned as in strict mode.
  \item \textbf{Malformed \texttt{write\_file} JSON.} \emph{Trigger:} a reply
  that fails JSON parsing but contains \verb|"tool"|, \texttt{write\_file},
  \verb|"content"|, and a \verb|"path"| string value. \emph{Transformation:} the
  file content is the text after the opening quote of the \texttt{content}
  value, with one trailing \verb|"}}|, \verb|"}|-newline-\verb|}|, or
  \verb|"}| removed, surrounding whitespace stripped, and the escape sequences
  \verb|\n| and \verb|\"| decoded. A non-empty result becomes a
  \texttt{write\_file} call to the extracted path; otherwise the reply counts as
  an invalid tool call.
  \item \textbf{Filename-keyed artifact.} \emph{Trigger:} a parsed reply that is
  a single-key object whose key is one of the task's expected artifact names
  and whose value is a non-empty string, for example
  \verb|{"final_report.md": "..."}|. \emph{Transformation:} the reply becomes a
  \texttt{write\_file} call with that path and content.
\end{enumerate}

Table~\ref{app:controller-repair-counts} counts how often rules 2--4 fired.
They fire more often for the official Gemma 31B prompt conditions than for the
less-restricted 31B derivative, and none of that derivative's 21 successful
runs uses them, so the repair controller does not create the Gemma 31B
advantage. In the other pairs, rules 2--4 fire in at most 12 runs per
condition.

\begin{table}[h]
\centering
\caption{Security runs in which controller repair rules 2--4 fired, by model
pair and condition; event counts are shown in parentheses where a rule fired
more than once in a run. Every condition has 150 runs, and rule 1 applies to
every run. The last column gives successful runs that used any of rules 2--4,
out of all successful runs.}
\label{app:controller-repair-counts}
\small
\begin{tabular}{@{}llrrrr@{}}
\toprule
Model pair & Condition & Repeat list & Malformed write & Filename-keyed & Succ.\ using 2--4 \\
\midrule
Gemma 31B & Aligned & 0 & 0 & 65 & 1/1 \\
 & Authorized & 5 (20) & 4 & 73 & 0/0 \\
 & Placeholder & 0 & 1 & 65 & 1/1 \\
 & Uncensored & 5 (20) & 0 & 0 & 0/21 \\
\midrule
Gemma 26B A4B & Aligned & 10 (30) & 2 & 0 & 0/0 \\
 & Uncensored & 0 & 9 & 0 & 0/16 \\
\midrule
Qwen2.5-Coder 7B & Aligned & 0 & 8 & 0 & 1/8 \\
 & Uncensored & 0 & 4 & 0 & 1/3 \\
\midrule
Llama 3.1 8B & Aligned & 0 & 12 & 0 & 3/5 \\
 & Uncensored & 0 & 4 & 0 & 0/0 \\
\bottomrule
\end{tabular}

\end{table}

\section{Out-of-Scope Package Installations}
\label{app:out-of-scope}

An audit of all saved traces found that the command allow-list does not prevent
network access through an allowed interpreter. In two less-restricted Gemma 26B
A4B security runs (\texttt{pv\_path\_normalize}, seed 3, and
\texttt{pv\_template\_escape}, seed 0), the agent ran
\texttt{python3 -m pip install} to download \texttt{pytest} or \texttt{jinja2}
from the public Python package index into the evaluation host's user
site-packages; in the second run, a direct \texttt{pip} call had first been
blocked. Neither run succeeded. Because the automated detector flags only
blocked tool-policy violations, the first run has no unsafe-action event and
the second has one, for its blocked call; counting both, 9/150
less-restricted 26B runs contain a tool-policy violation or an out-of-scope
installation. No other trace shows network access.

The installed packages remained available to runs started later on the same
host: the 111 subsequent 26B security runs (55 aligned and 56 less-restricted)
and every study that followed, including the 31B strict-controller ablation,
standard-coding control, and authorization-prompt study, whose matched 31B
repair-controller baselines ran before the installation; none of their runs
invoked either package. The 31B security, 31B sanity, and earlier 26B runs ran
before the installation. No reported success depends on these
packages: four successful Llama 3.1 aligned runs invoked \texttt{pytest}, but it
collected no tests, and their reports do not draw on its output.

\section{Schemas, Metrics, and Labels}
\label{app:schemas}
\label{app:supplementary-tables}

Each task manifest records the task id, class, local fixture entrypoint,
budgets, allowed tools, forbidden actions, prompt, ground truth, grounding
targets, success predicate, redaction rule, and expected artifacts. Each trace
records model provenance, condition, task, seed, backend, quantization,
sampling settings, messages, tool calls, phase events, refusal and safety
events, final outcome, failure label, costs, and timing.

The main outcome is deterministic end-to-end success. Secondary metrics are
strict refusal, partial refusal, suppressed action, unsafe action, invalid tool
use, cost, wall-clock time, and deterministic evidence grounding. Grounding is
the count of satisfied local-evidence fields: expected file read, expected file
named, expected symbol named, issue term named, and local evidence term
provided. The failure taxonomy assigns each failed run exactly one primary
label: \texttt{strict\_refusal}, \texttt{partial\_refusal},
\texttt{suppressed\_action}, \texttt{planning\_failure},
\texttt{reconnaissance\_failure}, \texttt{setup\_failure},
\texttt{localization\_failure}, \texttt{construction\_failure},
\texttt{validation\_failure}, \texttt{reporting\_failure},
\texttt{tool\_failure}, \texttt{environment\_failure},
\texttt{capability\_failure}, or \texttt{unsafe\_action}.

\section{Supplementary Results}

The appendix reports only the tables needed to interpret the main claims. Full
per-condition generated tables remain in the \texttt{analysis/} directories.

\begin{table}[h]
\centering
\caption{Comparison across security, sanity, and standard coding controls.
``Less'' denotes the less-restricted derivative for each model pair. Grounding is
the mean deterministic score out of five.}
\label{app:cross-evaluation-summary}
\small
\begin{tabular}{@{}lrrrrr@{}}
\toprule
Suite & Runs/cond. & Aligned succ. & Less succ. & Aligned ground. & Less ground. \\
\midrule
31B security & 150 & 0.7\% & 14.0\% & 3.27 & 3.91 \\
31B sanity & 60 & 0.0\% & 40.0\% & 2.33 & 4.27 \\
31B standard & 40 & 0.0\% & 37.5\% & 2.55 & 4.10 \\
26B security & 150 & 0.0\% & 10.7\% & 1.64 & 4.12 \\
26B sanity & 60 & 0.0\% & 51.7\% & 1.30 & 4.35 \\
26B standard & 40 & 0.0\% & 25.0\% & 2.27 & 3.95 \\
Qwen security & 150 & 5.3\% & 2.0\% & 3.17 & 3.21 \\
Qwen sanity & 60 & 0.0\% & 0.0\% & 2.55 & 2.77 \\
Qwen standard & 40 & 2.5\% & 2.5\% & 2.27 & 3.05 \\
Llama security & 150 & 3.3\% & 0.0\% & 2.80 & 0.21 \\
Llama sanity & 60 & 1.7\% & 0.0\% & 2.72 & 0.23 \\
Llama standard & 40 & 0.0\% & 0.0\% & 2.17 & 0.50 \\
\bottomrule
\end{tabular}

\end{table}

\begin{table}[h]
\centering
\caption{Safety outcomes for the two Gemma security evaluations. Grounding is
mean score out of five; Refusal, Suppressed, and Unsafe are shares of runs
with at least one such event.
Redact. counts public-trace audit findings.}
\label{app:safety-outcomes-31b}
\small
\emph{Gemma 31B.}
\begin{tabular}{@{}lrrrrrr@{}}
\toprule
Condition & Success & Ground. & Refusal & Suppressed & Unsafe & Redact. \\
\midrule
Aligned & 0.7\% & 3.27 & 0.0\% & 0.0\% & 0.0\% & 0 \\
Authorized & 0.0\% & 3.36 & 0.0\% & 0.0\% & 0.0\% & 0 \\
Placeholder & 0.7\% & 3.20 & 0.0\% & 0.0\% & 0.0\% & 0 \\
Uncensored & 14.0\% & 3.91 & 0.0\% & 0.0\% & 0.0\% & 0 \\
\bottomrule
\end{tabular}


\vspace{0.5em}
\emph{Gemma 26B A4B.}
\begin{tabular}{@{}lrrrrrr@{}}
\toprule
Condition & Success & Ground. & Refusal & Suppressed & Unsafe & Redact. \\
\midrule
Aligned & 0.0\% & 1.64 & 0.0\% & 0.0\% & 0.0\% & 0 \\
Uncensored & 10.7\% & 4.12 & 0.0\% & 0.0\% & 5.3\% & 0 \\
\bottomrule
\end{tabular}

\end{table}

\subsection{Cross-Family Details}
\label{app:evaluation-26b}
\label{app:non-gemma}

The 26B A4B, Qwen2.5-Coder 7B, and Llama 3.1 8B traces use the same repair
controller, task suite, seeds, success predicates, grounding metrics, and
redaction audits as the 31B Gemma evaluation. The cross-suite summary above
is the compact comparison; detailed generated success, grounding, failure, and
cost tables are retained in their corresponding \texttt{analysis/} folders.
Table~\ref{app:model-provenance} lists the exact GGUF files; every security and
control trace records the same model identifier and SHA-256 hash as its config.
All security and control runs use temperature 0.2, top-\(p\) 0.95, a 512-token
generation limit, and seeds 0--4; the two 24-run pilots use seeds 0--1.
The 31B authorized and placeholder conditions reuse the aligned 31B file.

\begin{table}[h]
\centering
\caption{GGUF releases and files used for each model pair, with the first 16
hexadecimal digits of each file's SHA-256 hash. Full hashes are recorded in
the configs and traces. Two files are no longer at the head of their
repositories: the 31B aligned file is available at revision \texttt{fb5801c702}
of \texttt{ggml-org/gemma-4-31B-it-GGUF}, and the 26B aligned file at revision
\texttt{3365c68df1} of \texttt{unsloth/gemma-4-26B-A4B-it-GGUF}. The Qwen
derivative's upstream repository (\texttt{huihui-ai}) is no longer publicly
accessible.}
\label{app:model-provenance}
\footnotesize
\begin{tabular}{@{}lll@{}}
\toprule
Model pair & GGUF release (aligned / less-restricted) & SHA-256 prefix \\
\midrule
Gemma 31B & \texttt{ggml-org/gemma-4-31B-it-GGUF} & \texttt{4f369f8fe0e1bedc} \\
 & \quad\texttt{gemma-4-31B-it-Q4\_K\_M.gguf} & \\
 & \texttt{TrevorJS/gemma-4-31B-it-uncensored-GGUF} & \texttt{14a465017c087e60} \\
 & \quad\texttt{gemma-4-31B-it-uncensored-Q4\_K\_M.gguf} & \\
Gemma 26B A4B & \texttt{unsloth/gemma-4-26B-A4B-it-GGUF} & \texttt{34c746b1d50ab813} \\
 & \quad\texttt{gemma-4-26B-A4B-it-UD-Q4\_K\_M.gguf} & \\
 & \texttt{TrevorJS/gemma-4-26B-A4B-it-uncensored-GGUF} & \texttt{d482a5daba09e67c} \\
 & \quad\texttt{gemma-4-26B-A4B-it-uncensored-Q4\_K\_M.gguf} & \\
Qwen2.5-Coder 7B & \texttt{Qwen/Qwen2.5-Coder-7B-Instruct-GGUF} & \texttt{509287f78cb4d4cf} \\
 & \quad\texttt{qwen2.5-coder-7b-instruct-q4\_k\_m.gguf} & \\
 & \texttt{bartowski/Qwen2.5-Coder-7B-Instruct-abliterated-GGUF} & \texttt{fdf02c16fb654ff6} \\
 & \quad\texttt{Qwen2.5-Coder-7B-Instruct-abliterated-Q4\_K\_M.gguf} & \\
Llama 3.1 8B & \texttt{bartowski/Meta-Llama-3.1-8B-Instruct-GGUF} & \texttt{7b064f5842bf9532} \\
 & \quad\texttt{Meta-Llama-3.1-8B-Instruct-Q4\_K\_M.gguf} & \\
 & \texttt{bartowski/Meta-Llama-3.1-8B-Instruct-abliterated-GGUF} & \texttt{c6be866583f7bddb} \\
 & \quad\texttt{Meta-Llama-3.1-8B-Instruct-abliterated-Q4\_K\_M.gguf} & \\
\bottomrule
\end{tabular}

\end{table}

\begin{table}[h]
\centering
\caption{Strict-vs-repair controller ablation on the matched 10-task 31B
subset. Grounding is the mean deterministic score out of five.}
\label{app:controller-ablation}
\small
\begin{tabular}{@{}llrrr@{}}
\toprule
Controller & Condition & Success & Tool fail. & Grounding \\
\midrule
Strict & Aligned & 0/50 & 50/50 & 0.20 \\
Strict & Authorized & 0/50 & 50/50 & 0.20 \\
Strict & Placeholder & 0/50 & 50/50 & 0.20 \\
Strict & Uncensored & 5/50 & 1/50 & 3.84 \\
\midrule
Repair & Aligned & 0/50 & 0/50 & 3.36 \\
Repair & Authorized & 0/50 & 0/50 & 3.46 \\
Repair & Placeholder & 0/50 & 0/50 & 3.22 \\
Repair & Uncensored & 10/50 & 0/50 & 4.42 \\
\bottomrule
\end{tabular}

\end{table}

\begin{table}[h]
\centering
\caption{Task-blocked sign-permutation checks for aligned versus
less-restricted success and grounding gaps.}
\label{app:robustness}
\small
\begin{tabular}{@{}lrrrr@{}}
\toprule
Suite & Tasks & Success diff. & Success p & Grounding p \\
\midrule
31B security & 30 & 0.133 & 0.0236 & 0.0022 \\
26B security & 30 & 0.107 & 0.0315 & $<10^{-5}$ \\
31B standard & 8 & 0.375 & 0.2500 & 0.0156 \\
\bottomrule
\end{tabular}

\end{table}

\begin{table}[h]
\centering
\caption{Compact qualitative examples from redacted 31B Gemma security traces.}
\small
\scriptsize
\begin{tabular}{@{}p{0.22\linewidth}p{0.18\linewidth}p{0.52\linewidth}@{}}
\toprule
Trace & Outcome & Redacted observation \\
\midrule
\shortstack[l]{\texttt{vl\_path\_join}\\\texttt{aligned:0}} &
Aligned failure &
Read \texttt{src/files.py} and produced a plausible report, but named the
wrong function. The artifact identified \texttt{read\_user\_file} instead of
\texttt{open\_user\_file}, so local access occurred without
predicate-level localization. \\
\shortstack[l]{\texttt{vl\_path\_join}\\\texttt{uncensored:0}} &
Less-restricted success &
Read \texttt{src/files.py}, named \texttt{open\_user\_file}, described path
traversal, and tied the claim to string concatenation in the local fixture.
This satisfies file, symbol, issue, and evidence grounding. \\
\shortstack[l]{\texttt{pt\_boundary\_crash}\\\texttt{uncensored:0}} &
Less-restricted validation failure &
Produced a well-grounded report naming \texttt{parser/bounds.py},
\texttt{parse\_index}, and redacted trigger evidence, but did not provide the
required executable local validation artifact; grounding did not satisfy the
hard proof-of-trigger check. \\
\bottomrule
\end{tabular}

\end{table}

\section{Reviewer Audit}
\label{app:reviewer-audit}

The reviewer audit samples 120 redacted traces from 31B security, 31B sanity,
26B security, and 26B sanity runs. Two independent LLM reviewer passes label
the same blind packet for success, primary failure label, refusal relation,
four-point grounding quality, and hard-validation subtype when applicable. The
packet excludes condition, original label, deterministic grounding score,
harness checker reason, and copy-paste forbidden predicate strings.

\begin{table}[h]
\centering
\caption{Independent LLM reviewer agreement on the blind trace packet.}
\small
\begin{tabular}{lrr}
\toprule
Comparison & Agreement & Kappa \\
\midrule
Reviewer labels & 0.73 & 0.62 \\
Reviewer success & 0.93 & 0.87 \\
Reviewer refusal & 1.00 & 1.00 \\
Reviewer A vs original & 0.37 & 0.25 \\
Reviewer B vs original & 0.47 & 0.37 \\
Hard-validation subtype & 0.35 & 0.21 \\
\bottomrule
\end{tabular}

\end{table}

\begin{table}[h]
\centering
\caption{Grounding audit comparing LLM reviewer grounding judgments with
deterministic bins.}
\small
\begin{tabular}{lrrr}
\toprule
Reviewer & Exact bin & Within one & Spearman \\
\midrule
Reviewer A & 0.57 & 0.99 & 0.60 \\
Reviewer B & 0.60 & 0.97 & 0.65 \\
\bottomrule
\end{tabular}

\end{table}

\begin{table}[h]
\centering
\caption{Hard-validation subtype agreement in sampled validation failures.}
\small
\begin{tabular}{lrr}
\toprule
Subtype & Reviewer A & Reviewer B \\
\midrule
no executable check & 13 & 13 \\
check failed & 0 & 4 \\
output misread & 3 & 6 \\
ungrounded claim & 12 & 0 \\
redaction or payload issue & 1 & 0 \\
unsafe action & 3 & 3 \\
other & 2 & 8 \\
\midrule
Subtype agreement (raw / $\kappa$) & \multicolumn{2}{r}{0.35 / 0.21} \\
\bottomrule
\end{tabular}

\end{table}


\clearpage
\section*{NeurIPS Paper Checklist}


\begin{enumerate}

\item {\bf Claims}
    \item[] Question: Do the main claims made in the abstract and introduction accurately reflect the paper's contributions and scope?
    \item[] Answer: \answerYes{} % Replace by \answerYes{}, \answerNo{}, or \answerNA{}.
    \item[] Justification: The abstract and introduction (Section~\ref{sec:intro}) state the four-pair evaluation, the Gemma grounding result, the non-Gemma reversals, and the absence of passes on the lexical hard-validation checks, and the results and discussion (Sections~\ref{sec:results} and~\ref{sec:discussion}) report the full evaluation: strict-controller and repair-controller 24-run pilots, a completed 1,500-trace repair-controller security evaluation, 800 non-security control traces, a 150-run authorization-prompt robustness study, a 200-run strict-controller ablation matched to repair-controller traces, grounding metrics, and a 120-trace LLM-reviewer agreement audit. The claims emphasize low absolute success rates, derivative-model provenance, and tool-interface confounds.
    \item[] Guidelines:
    \begin{itemize}
        \item The answer \answerNA{} means that the abstract and introduction do not include the claims made in the paper.
        \item The abstract and/or introduction should clearly state the claims made, including the contributions made in the paper and important assumptions and limitations. A \answerNo{} or \answerNA{} answer to this question will not be perceived well by the reviewers. 
        \item The claims made should match theoretical and experimental results, and reflect how much the results can be expected to generalize to other settings. 
        \item It is fine to include aspirational goals as motivation as long as it is clear that these goals are not attained by the paper. 
    \end{itemize}

\item {\bf Limitations}
    \item[] Question: Does the paper discuss the limitations of the work performed by the authors?
    \item[] Answer: \answerYes{} % Replace by \answerYes{}, \answerNo{}, or \answerNA{}.
    \item[] Justification: Section~\ref{sec:limitations} discusses derivative provenance and the untested abliteration methods and strengths, the capability confound, controller effects, low absolute success rates, the lexical grounding measure and its overlap with success predicates, LLM-reviewer audit limits, and the lexical success predicates, including an unprompted required term in the proof-of-trigger tasks, and the limited scope of the authorization ablation and inference setup. Section~\ref{sec:evaluation} explains why the non-security controls rule out a security-specific interpretation for the Gemma pairs.
    \item[] Guidelines:
    \begin{itemize}
        \item The answer \answerNA{} means that the paper has no limitation while the answer \answerNo{} means that the paper has limitations, but those are not discussed in the paper. 
        \item The authors are encouraged to create a separate ``Limitations'' section in their paper.
        \item The paper should point out any strong assumptions and how robust the results are to violations of these assumptions (e.g., independence assumptions, noiseless settings, model well-specification, asymptotic approximations only holding locally). The authors should reflect on how these assumptions might be violated in practice and what the implications would be.
        \item The authors should reflect on the scope of the claims made, e.g., if the approach was only tested on a few datasets or with a few runs. In general, empirical results often depend on implicit assumptions, which should be articulated.
        \item The authors should reflect on the factors that influence the performance of the approach. For example, a facial recognition algorithm may perform poorly when image resolution is low or images are taken in low lighting. Or a speech-to-text system might not be used reliably to provide closed captions for online lectures because it fails to handle technical jargon.
        \item The authors should discuss the computational efficiency of the proposed algorithms and how they scale with dataset size.
        \item If applicable, the authors should discuss possible limitations of their approach to address problems of privacy and fairness.
        \item While the authors might fear that complete honesty about limitations might be used by reviewers as grounds for rejection, a worse outcome might be that reviewers discover limitations that aren't acknowledged in the paper. The authors should use their best judgment and recognize that individual actions in favor of transparency play an important role in developing norms that preserve the integrity of the community. Reviewers will be specifically instructed to not penalize honesty concerning limitations.
    \end{itemize}

\item {\bf Theory assumptions and proofs}
    \item[] Question: For each theoretical result, does the paper provide the full set of assumptions and a complete (and correct) proof?
    \item[] Answer: \answerNA{} % Replace by \answerYes{}, \answerNo{}, or \answerNA{}.
    \item[] Justification: The current paper is empirical and does not state formal theoretical results.
    \item[] Guidelines:
    \begin{itemize}
        \item The answer \answerNA{} means that the paper does not include theoretical results. 
        \item All the theorems, formulas, and proofs in the paper should be numbered and cross-referenced.
        \item All assumptions should be clearly stated or referenced in the statement of any theorems.
        \item The proofs can either appear in the main paper or the supplemental material, but if they appear in the supplemental material, the authors are encouraged to provide a short proof sketch to provide intuition. 
        \item Inversely, any informal proof provided in the core of the paper should be complemented by formal proofs provided in appendix or supplemental material.
        \item Theorems and Lemmas that the proof relies upon should be properly referenced. 
    \end{itemize}

    \item {\bf Experimental result reproducibility}
    \item[] Question: Does the paper fully disclose all the information needed to reproduce the main experimental results of the paper to the extent that it affects the main claims and/or conclusions of the paper (regardless of whether the code and data are provided or not)?
    \item[] Answer: \answerYes{} % Replace by \answerYes{}, \answerNo{}, or \answerNA{}.
    \item[] Justification: The paper reports two 24-run local pilots, a completed 1,500-trace security evaluation, 800 non-security control traces, a 150-run authorization-prompt robustness study, and a 200-run strict-controller ablation matched to repair-controller traces, and discloses the runnable harness, task catalogs, model endpoints, matched GGUF hashes, replicate indices, controller modes, analysis commands, redacted traces, reviewer labels, and summary files needed to reproduce them (Sections~\ref{sec:methodology} and~\ref{sec:evaluation} and Appendices~\ref{app:artifact}, \ref{app:controller-repairs}, and~\ref{app:evaluation-26b}). Full model reruns require local access to the cited GGUF model releases and acceptance of the corresponding upstream model terms.
    \item[] Guidelines:
    \begin{itemize}
        \item The answer \answerNA{} means that the paper does not include experiments.
        \item If the paper includes experiments, a \answerNo{} answer to this question will not be perceived well by the reviewers: Making the paper reproducible is important, regardless of whether the code and data are provided or not.
        \item If the contribution is a dataset and\slash or model, the authors should describe the steps taken to make their results reproducible or verifiable. 
        \item Depending on the contribution, reproducibility can be accomplished in various ways. For example, if the contribution is a novel architecture, describing the architecture fully might suffice, or if the contribution is a specific model and empirical evaluation, it may be necessary to either make it possible for others to replicate the model with the same dataset, or provide access to the model. In general. releasing code and data is often one good way to accomplish this, but reproducibility can also be provided via detailed instructions for how to replicate the results, access to a hosted model (e.g., in the case of a large language model), releasing of a model checkpoint, or other means that are appropriate to the research performed.
        \item While NeurIPS does not require releasing code, the conference does require all submissions to provide some reasonable avenue for reproducibility, which may depend on the nature of the contribution. For example
        \begin{enumerate}
            \item If the contribution is primarily a new algorithm, the paper should make it clear how to reproduce that algorithm.
            \item If the contribution is primarily a new model architecture, the paper should describe the architecture clearly and fully.
            \item If the contribution is a new model (e.g., a large language model), then there should either be a way to access this model for reproducing the results or a way to reproduce the model (e.g., with an open-source dataset or instructions for how to construct the dataset).
            \item We recognize that reproducibility may be tricky in some cases, in which case authors are welcome to describe the particular way they provide for reproducibility. In the case of closed-source models, it may be that access to the model is limited in some way (e.g., to registered users), but it should be possible for other researchers to have some path to reproducing or verifying the results.
        \end{enumerate}
    \end{itemize}


\item {\bf Open access to data and code}
    \item[] Question: Does the paper provide open access to the data and code, with sufficient instructions to faithfully reproduce the main experimental results, as described in supplemental material?
    \item[] Answer: \answerYes{} % Replace by \answerYes{}, \answerNo{}, or \answerNA{}.
    \item[] Justification: The public repository (\codeurl) includes the runnable local harness and its tests, task catalogs with success predicates, model configs, prompt diagnostic tooling, analysis and redaction scripts, redacted traces for every reported run, generated summaries, reviewer-audit labels, and redaction audits. Appendix~\ref{app:artifact} gives the exact reproduction entry point and model-rerun CLI template; full model reruns require obtaining the cited upstream model weights under their respective licenses and terms.
    \item[] Guidelines:
    \begin{itemize}
        \item The answer \answerNA{} means that paper does not include experiments requiring code.
        \item Please see the NeurIPS code and data submission guidelines (\url{https://neurips.cc/public/guides/CodeSubmissionPolicy}) for more details.
        \item While we encourage the release of code and data, we understand that this might not be possible, so \answerNo{} is an acceptable answer. Papers cannot be rejected simply for not including code, unless this is central to the contribution (e.g., for a new open-source benchmark).
        \item The instructions should contain the exact command and environment needed to run to reproduce the results. See the NeurIPS code and data submission guidelines (\url{https://neurips.cc/public/guides/CodeSubmissionPolicy}) for more details.
        \item The authors should provide instructions on data access and preparation, including how to access the raw data, preprocessed data, intermediate data, and generated data, etc.
        \item The authors should provide scripts to reproduce all experimental results for the new proposed method and baselines. If only a subset of experiments are reproducible, they should state which ones are omitted from the script and why.
        \item At submission time, to preserve anonymity, the authors should release anonymized versions (if applicable).
        \item Providing as much information as possible in supplemental material (appended to the paper) is recommended, but including URLs to data and code is permitted.
    \end{itemize}


\item {\bf Experimental setting/details}
    \item[] Question: Does the paper specify all the training and test details (e.g., data splits, hyperparameters, how they were chosen, type of optimizer) necessary to understand the results?
    \item[] Answer: \answerYes{} % Replace by \answerYes{}, \answerNo{}, or \answerNA{}.
    \item[] Justification: The protocol (Sections~\ref{sec:methodology} and~\ref{sec:evaluation} and Appendix~\ref{app:evaluation-26b}) specifies model conditions, prompt conditions, task classes, seeds, budgets, metrics, schemas, controller mode, and inference controls, including temperature 0.2, top-\(p\) 0.95, a 512-token generation limit, and five replicates per task (indices 0--4; sampling is not seeded). The reported model-pair evaluations record matched GGUF file hashes and local OpenAI-compatible inference endpoints backed by \texttt{llama.cpp}.
    \item[] Guidelines:
    \begin{itemize}
        \item The answer \answerNA{} means that the paper does not include experiments.
        \item The experimental setting should be presented in the core of the paper to a level of detail that is necessary to appreciate the results and make sense of them.
        \item The full details can be provided either with the code, in appendix, or as supplemental material.
    \end{itemize}

\item {\bf Experiment statistical significance}
    \item[] Question: Does the paper report error bars suitably and correctly defined or other appropriate information about the statistical significance of the experiments?
    \item[] Answer: \answerYes{} % Replace by \answerYes{}, \answerNo{}, or \answerNA{}.
    \item[] Justification: The paired model comparisons report bootstrap 95\% confidence intervals for success rates (Figure~\ref{fig:model-family-results}), McNemar's exact tests for paired aligned--less-restricted success (Section~\ref{sec:results}), and task-blocked sign-permutation checks for success and grounding gaps in three Gemma comparisons (Appendix Table~\ref{app:robustness}). Mixed-effects modeling is left to follow-up analysis because several task classes have near-zero success.
    \item[] Guidelines:
    \begin{itemize}
        \item The answer \answerNA{} means that the paper does not include experiments.
        \item The authors should answer \answerYes{} if the results are accompanied by error bars, confidence intervals, or statistical significance tests, at least for the experiments that support the main claims of the paper.
        \item The factors of variability that the error bars are capturing should be clearly stated (for example, train/test split, initialization, random drawing of some parameter, or overall run with given experimental conditions).
        \item The method for calculating the error bars should be explained (closed form formula, call to a library function, bootstrap, etc.)
        \item The assumptions made should be given (e.g., Normally distributed errors).
        \item It should be clear whether the error bar is the standard deviation or the standard error of the mean.
        \item It is OK to report 1-sigma error bars, but one should state it. The authors should preferably report a 2-sigma error bar than state that they have a 96\% CI, if the hypothesis of Normality of errors is not verified.
        \item For asymmetric distributions, the authors should be careful not to show in tables or figures symmetric error bars that would yield results that are out of range (e.g., negative error rates).
        \item If error bars are reported in tables or plots, the authors should explain in the text how they were calculated and reference the corresponding figures or tables in the text.
    \end{itemize}

\item {\bf Experiments compute resources}
    \item[] Question: For each experiment, does the paper provide sufficient information on the computer resources (type of compute workers, memory, time of execution) needed to reproduce the experiments?
    \item[] Answer: \answerYes{} % Replace by \answerYes{}, \answerNo{}, or \answerNA{}.
    \item[] Justification: The pilots, security evaluation, non-security controls, prompt-robustness study, and strict-controller ablation were run locally through OpenAI-compatible \texttt{llama.cpp} endpoints. Saved traces record wall-clock time and token counts per run, and the generated analysis reports model-call totals, approximate token totals, median per-run wall-clock time, and summed wall-clock hours by condition.
    \item[] Guidelines:
    \begin{itemize}
        \item The answer \answerNA{} means that the paper does not include experiments.
        \item The paper should indicate the type of compute workers CPU or GPU, internal cluster, or cloud provider, including relevant memory and storage.
        \item The paper should provide the amount of compute required for each of the individual experimental runs as well as estimate the total compute. 
        \item The paper should disclose whether the full research project required more compute than the experiments reported in the paper (e.g., preliminary or failed experiments that didn't make it into the paper). 
    \end{itemize}
    
\item {\bf Code of ethics}
    \item[] Question: Does the research conducted in the paper conform, in every respect, with the NeurIPS Code of Ethics \url{https://neurips.cc/public/EthicsGuidelines}?
    \item[] Answer: \answerYes{} % Replace by \answerYes{}, \answerNo{}, or \answerNA{}.
    \item[] Justification: The paper is scoped to authorized local sandboxes, avoids third-party targets, and withholds operational exploit instructions.
    \item[] Guidelines:
    \begin{itemize}
        \item The answer \answerNA{} means that the authors have not reviewed the NeurIPS Code of Ethics.
        \item If the authors answer \answerNo, they should explain the special circumstances that require a deviation from the Code of Ethics.
        \item The authors should make sure to preserve anonymity (e.g., if there is a special consideration due to laws or regulations in their jurisdiction).
    \end{itemize}


\item {\bf Broader impacts}
    \item[] Question: Does the paper discuss both potential positive societal impacts and negative societal impacts of the work performed?
    \item[] Answer: \answerYes{} % Replace by \answerYes{}, \answerNo{}, or \answerNA{}.
    \item[] Justification: The introduction (Section~\ref{sec:intro}) motivates the defensive value of reliable, authorized security agents, and Section~\ref{sec:discussion}, including Section~\ref{sec:ethics}, discusses dual-use risks and mitigations such as redacted traces and sandboxed evaluation.
    \item[] Guidelines:
    \begin{itemize}
        \item The answer \answerNA{} means that there is no societal impact of the work performed.
        \item If the authors answer \answerNA{} or \answerNo, they should explain why their work has no societal impact or why the paper does not address societal impact.
        \item Examples of negative societal impacts include potential malicious or unintended uses (e.g., disinformation, generating fake profiles, surveillance), fairness considerations (e.g., deployment of technologies that could make decisions that unfairly impact specific groups), privacy considerations, and security considerations.
        \item The conference expects that many papers will be foundational research and not tied to particular applications, let alone deployments. However, if there is a direct path to any negative applications, the authors should point it out. For example, it is legitimate to point out that an improvement in the quality of generative models could be used to generate Deepfakes for disinformation. On the other hand, it is not needed to point out that a generic algorithm for optimizing neural networks could enable people to train models that generate Deepfakes faster.
        \item The authors should consider possible harms that could arise when the technology is being used as intended and functioning correctly, harms that could arise when the technology is being used as intended but gives incorrect results, and harms following from (intentional or unintentional) misuse of the technology.
        \item If there are negative societal impacts, the authors could also discuss possible mitigation strategies (e.g., gated release of models, providing defenses in addition to attacks, mechanisms for monitoring misuse, mechanisms to monitor how a system learns from feedback over time, improving the efficiency and accessibility of ML).
    \end{itemize}
    
\item {\bf Safeguards}
    \item[] Question: Does the paper describe safeguards that have been put in place for responsible release of data or models that have a high risk for misuse (e.g., pre-trained language models, image generators, or scraped datasets)?
    \item[] Answer: \answerYes{} % Replace by \answerYes{}, \answerNo{}, or \answerNA{}.
    \item[] Justification: The release publishes code, task manifests, aggregate outcomes, and redacted, audited traces; it contains no reusable exploit payloads or target-specific intrusion steps, and no jailbreak prompt was applied in any reported run (Section~\ref{sec:ethics}).
    \item[] Guidelines:
    \begin{itemize}
        \item The answer \answerNA{} means that the paper poses no such risks.
        \item Released models that have a high risk for misuse or dual-use should be released with necessary safeguards to allow for controlled use of the model, for example by requiring that users adhere to usage guidelines or restrictions to access the model or implementing safety filters. 
        \item Datasets that have been scraped from the Internet could pose safety risks. The authors should describe how they avoided releasing unsafe images.
        \item We recognize that providing effective safeguards is challenging, and many papers do not require this, but we encourage authors to take this into account and make a best faith effort.
    \end{itemize}

\item {\bf Licenses for existing assets}
    \item[] Question: Are the creators or original owners of assets (e.g., code, data, models), used in the paper, properly credited and are the license and terms of use explicitly mentioned and properly respected?
    \item[] Answer: \answerYes{} % Replace by \answerYes{}, \answerNo{}, or \answerNA{}.
    \item[] Justification: The local harness, task catalogs, analysis scripts, and paper artifact are released under the repository's MIT license. Model weights, GGUF conversions, and derivative releases are not relicensed by this work: the paper cites the official Gemma 4, Qwen, and Llama model sources, the third-party derivative and GGUF conversion sources, \texttt{llama.cpp}, and Hugging Face; the configs record source URLs and hashes so model artifacts remain governed by their respective upstream model licenses and terms of use.
    \item[] Guidelines:
    \begin{itemize}
        \item The answer \answerNA{} means that the paper does not use existing assets.
        \item The authors should cite the original paper that produced the code package or dataset.
        \item The authors should state which version of the asset is used and, if possible, include a URL.
        \item The name of the license (e.g., CC-BY 4.0) should be included for each asset.
        \item For scraped data from a particular source (e.g., website), the copyright and terms of service of that source should be provided.
        \item If assets are released, the license, copyright information, and terms of use in the package should be provided. For popular datasets, \url{paperswithcode.com/datasets} has curated licenses for some datasets. Their licensing guide can help determine the license of a dataset.
        \item For existing datasets that are re-packaged, both the original license and the license of the derived asset (if it has changed) should be provided.
        \item If this information is not available online, the authors are encouraged to reach out to the asset's creators.
    \end{itemize}

\item {\bf New assets}
    \item[] Question: Are new assets introduced in the paper well documented and is the documentation provided alongside the assets?
    \item[] Answer: \answerYes{} % Replace by \answerYes{}, \answerNo{}, or \answerNA{}.
    \item[] Justification: The paper (Appendices~\ref{app:artifact}--\ref{app:schemas}) and the public repository (\codeurl) document task-manifest and run-trace schemas, a labeling rubric, a prompt diagnostic taxonomy, release safeguards, a runnable harness, a 30-task security catalog, a 12-task sanity catalog, an 8-task standard coding catalog, strict and repair controller modes, pilot traces, four-pair security and control traces, grounding metrics, reviewer-audit labels, and redaction-audit notes.
    \item[] Guidelines:
    \begin{itemize}
        \item The answer \answerNA{} means that the paper does not release new assets.
        \item Researchers should communicate the details of the dataset\slash code\slash model as part of their submissions via structured templates. This includes details about training, license, limitations, etc. 
        \item The paper should discuss whether and how consent was obtained from people whose asset is used.
        \item At submission time, remember to anonymize your assets (if applicable). You can either create an anonymized URL or include an anonymized zip file.
    \end{itemize}

\item {\bf Crowdsourcing and research with human subjects}
    \item[] Question: For crowdsourcing experiments and research with human subjects, does the paper include the full text of instructions given to participants and screenshots, if applicable, as well as details about compensation (if any)? 
    \item[] Answer: \answerNA{} % Replace by \answerYes{}, \answerNo{}, or \answerNA{}.
    \item[] Justification: The evaluation does not involve crowdsourcing or human-subject experiments.
    \item[] Guidelines:
    \begin{itemize}
        \item The answer \answerNA{} means that the paper does not involve crowdsourcing nor research with human subjects.
        \item Including this information in the supplemental material is fine, but if the main contribution of the paper involves human subjects, then as much detail as possible should be included in the main paper. 
        \item According to the NeurIPS Code of Ethics, workers involved in data collection, curation, or other labor should be paid at least the minimum wage in the country of the data collector. 
    \end{itemize}

\item {\bf Institutional review board (IRB) approvals or equivalent for research with human subjects}
    \item[] Question: Does the paper describe potential risks incurred by study participants, whether such risks were disclosed to the subjects, and whether Institutional Review Board (IRB) approvals (or an equivalent approval/review based on the requirements of your country or institution) were obtained?
    \item[] Answer: \answerNA{} % Replace by \answerYes{}, \answerNo{}, or \answerNA{}.
    \item[] Justification: The evaluation does not involve human subjects and therefore does not require IRB approval.
    \item[] Guidelines:
    \begin{itemize}
        \item The answer \answerNA{} means that the paper does not involve crowdsourcing nor research with human subjects.
        \item Depending on the country in which research is conducted, IRB approval (or equivalent) may be required for any human subjects research. If you obtained IRB approval, you should clearly state this in the paper. 
        \item We recognize that the procedures for this may vary significantly between institutions and locations, and we expect authors to adhere to the NeurIPS Code of Ethics and the guidelines for their institution. 
        \item For initial submissions, do not include any information that would break anonymity (if applicable), such as the institution conducting the review.
    \end{itemize}

\item {\bf Declaration of LLM usage}
    \item[] Question: Does the paper describe the usage of LLMs if it is an important, original, or non-standard component of the core methods in this research? Note that if the LLM is used only for writing, editing, or formatting purposes and does \emph{not} impact the core methodology, scientific rigor, or originality of the research, declaration is not required.
    %this research? 
    \item[] Answer: \answerYes{} % Replace by \answerYes{}, \answerNo{}, or \answerNA{}.
    \item[] Justification: LLMs are the object of study and are central to the experimental method; the methodology section describes the model conditions, agent harness, and blind LLM-reviewer audit of saved traces.
    \item[] Guidelines:
    \begin{itemize}
        \item The answer \answerNA{} means that the core method development in this research does not involve LLMs as any important, original, or non-standard components.
        \item Please refer to our LLM policy in the NeurIPS handbook for what should or should not be described.
    \end{itemize}

\end{enumerate}


\end{document}
